\chapter*{\uppercase{ABSTRACT}}
\addcontentsline{toc}{section}{\bfseries \uppercase{Abstract}}

Large-language-model applications are evolving from single-turn chat interfaces into systems that plan, use tools, preserve context, and coordinate multiple specialised agents. However, application developers commonly rebuild scheduling, memory, event propagation, execution monitoring, and failure handling for each project. This report presents AgentOS, a hybrid multi-agent runtime environment that provides these operating-system-style services around autonomous AI applications. The implemented system combines a Next.js user interface, a FastAPI gateway, a Redis-backed event path, MongoDB task history, ChromaDB semantic retrieval, structured LLM planning, dependency-aware directed acyclic graph execution, semantic context paging, a global priority-aging scheduler, token and concurrency admission control, and a sandboxed Python code runner with bounded self-healing retries. The frontend exposes the generated plan graph and live scheduler and execution activity through WebSocket updates. The report distinguishes implemented facilities from proposed components such as complete role-based access control, distributed locking, hybrid model routing, and durable coroutine restoration after a kernel restart. Verification is based on repository compile checks, TypeScript checks, and focused smoke tests rather than unsubstantiated performance claims. AgentOS therefore provides a modular foundation for studying reliable, observable, and resource-aware multi-agent execution.

\textbf{Keywords:} multi-agent runtime, semantic planning, DAG execution, context paging, priority scheduling.
